\subsection{构造 Lefschetz 纤维化}\label{sec:16e}
定义集族 $\symbfit V$ 上的 Hurwitz 移动时，并未预先要求它们来自消失圈。结果表明，这并无区别：任意这样的集族都可实现为某个纤维化的特选基。现在将此说得更准确些。

\clearpage
\StatementNumber{8}\begin{lemma}\label{st:16.8}
令 $V$ 为带角正合辛流形 $M$ 中带标架的正合 Lagrangian 球面。取圆盘 $S=D$，及从某点 $\gamma(0)\in\partial S$ 到 $\gamma(1)=0$ 的嵌入路径。则存在正合 Lefschetz 纤维化 $\pi:E\to S$，满足 $\operatorname{Critv}(\pi)=\{0\}$，并带有认同 $E_{\gamma(0)}=M$，使 $\gamma$ 的消失圈在此认同下对应于 $V$，且与标架相容。此外，该纤维化在 $\partial^hE$ 附近平凡。
\end{lemma}

构造很简单。从\stref{例题}{15.4}的模型出发，其纤维为球面余切丛的紧部分 $B_s^*S^n$，其中 $s$ 充分小。如 Dehn 扭转的定义，再用 Lagrangian 管状邻域定理找到延伸标架 $v:S^n\to V$ 的嵌入 $\psi:B_s^*S^n\to M$。借此把局部模型与平凡纤维化 $D\times(M\setminus V)$ 粘合。纤维上的辛形式一致，但在总空间上，还须用适当的截断函数把它们拼接。所需的只是，在相应 Lefschetz 顶针\eref{16.4}的邻域中保留局部模型的二形式；随后，所有消失路径都会给出所需的 Lagrangian 球面。

一旦能够构造具有单个指定消失圈的 Lefschetz 纤维化，就很容易将它们与非奇异部分粘合，得到一般结果：

\StatementNumber{9}\begin{lemma}\label{st:16.9}
令 $\symbfit V=(V_1,\ldots,V_m)$ 为带角正合辛流形 $M$ 中带标架正合 Lagrangian 球面的集族。在圆盘 $S=D$ 上选择基点 $*$，必要时还选择基路径 $\chi$，并选择消失路径特选基 $(\gamma_1,\ldots,\gamma_m)$。则存在正合 Lefschetz 纤维化 $\pi:E\to S=D$，其临界值恰是端点 $\gamma_1(1),\ldots,\gamma_m(1)$；它带有认同 $E_*=M$，使带标架消失圈 $V_{\gamma_k}$ 在此认同下对应于 $V_k$。此外，该纤维化在 $\partial^hE$ 附近平凡。
\end{lemma}

\stref{引理}{16.8}有一个唯一性方面的对应结果。具体地，假设在 $D$ 上有一个 Lefschetz 纤维化，只有一个临界点，位于 $0$ 上方的纤维中，
\SourcePage{237}
且在 $\partial^hE$ 附近平凡。最后这个要求是技术性的，重要性很有限；它源于 $\partial E$ 应弱 $I_E$ 凸这一条件，而该条件在辛联络变形下表现得不太方便。则在沿 $\partial^hE$ 行为良好的变形意义下，该纤维化完全由 $1$ 上方的纤维及其中的消失圈确定；消失圈视为同痕意义下的带标架 Lagrangian 球面。在\stref{引理}{16.9}更复杂的情形中也有相应结论。

\subsection{分次问题}\label{sec:16f}
令 $M^{2n}$ 为带角正合辛流形，带二次复体积形式 $\eta_M^2$。若 $V\subset M$ 为可赋予分次的带标架正合 Lagrangian 球面，并不需要具体选择某个分次，则 Dehn 扭转 $\tau_V$ 自然地成为\secref{12i}意义下的分次辛自同构；实值分次函数 $\gamma_{\tau_V}^\#$ 由下述条件确定：在 $V$ 的某个邻域之外恒为零，而 $\tau_V$ 在该处为恒等。证明并不困难：上同调检查表明，维数 $n>1$ 时不存在障碍，余下的情形则是初等的。\znote{229}{“分次函数” $\gamma^\#_{\tau_V}$ 应与\secref{12i}的定义对照读：$\tau_V$ 在 $V$ 之外为恒等，故 $\gamma^\#_{\tau_V}$ 支撑在 $V$ 附近，它不是 $\tau_V$ 的“像”，而是记录 $\tau_V$ 把分次 Lagrangian 的相位移动了多少的实值函数。$n>1$ 时 $V$ 的余维至少为 $2$，绕过 $V$ 的回路不会被迫穿过 $V$，相位提升的障碍自动消失；$n=1$ 时 $V$ 是余维 $1$ 的圈，$M\setminus V$ 被真正切开，两侧的相位提升可以相差整数，故须单独核对。$\eta_{E/S}^2=\eta_E^2/\eta_S^2$ 只是“相对”二字的含义：底空间的贡献由 $\eta_S^2$ 除掉，留在纤维上的是相对量。[来源：原书 §(16f)，印刷页 229。]}

另一方面，考虑配备相对二次复体积形式 $\eta_{E/S}^2$ 的 Lefschetz 纤维化。它在光滑纤维上诱导体积形式，关于这些形式，任意消失圈均可赋予分次，但没有自然选定的分次。可如前论证：除最低维情形外，其余情况都平凡，而最低维情形可直接解决。不过，还有一个更优美的论证，与证明消失圈正合性的论证类似：令 $\gamma$ 为消失路径。在 $\pi^{-1}(\gamma)\subset E$ 的某个邻域中，写成 $\eta_{E/S}^2=\eta_E^2/\eta_S^2$；事实上，由假设 $TS$ 平凡，可在整个空间中这样写，但此处不需要。由于 $\Delta_\gamma$ 可缩，与 $\eta_E^2$ 关联的平方相位函数
\[
\alpha_{\Delta_\gamma}:\Delta_\gamma\longrightarrow S^1
\]
当然可以提升到 $\R$。该函数在边界上的限制与 $V_\gamma$ 的平方相位函数只差一个依赖于 $\gamma'(0)$ 的常数，故得到所需事实。最后，考察上面概述的\eref{16.6}证明可见：若为单值变换赋予\eref{15.8}定义的分次，并按刚才描述的方式为 Dehn 扭转赋予分次，则那里构造的同痕提升到分次辛自同构群 $\operatorname{Aut}^{\mathrm{gr}}(M)$。

最后，若给定 $M$、$\eta_M^2$，以及一族均可赋予分次的带标架正合 Lagrangian 球面，则按\stref{引理}{16.9}构造的相应 Lefschetz 纤维化可赋予相对二次复体积形式，且相对于 $E_*=M$ 在同伦意义下唯一。同样，无须具体选择分次，只要它存在即可。

\StatementNumber{10}\begin{note}\label{st:16.10}
高维 Dehn 扭转为辛的事实，由 Arnol'd\bcite{3}发现。相应地，消失圈为 Lagrangian 这一事实
\SourcePage{238}
由 Donaldson 向作者指出，尽管此前也很可能已有人知道。至于 Picard–Lefschetz 定理，它在拓扑层次，即将单值变换视为微分同胚的层次，是经典结果。众多可用参考文献中的两项是\bcite{92}、\bcite{5}。辛版本见\bcite{131}命题 1.15。消失圈的特选基也是经典概念，在奇点理论考察的局部情形中尤其如此\bcite{5}、\bcite{4}。在辛层次上可从这种基重建总空间的观察，则较为近期，见本章开头引用的论文。Dehn 扭转的分次在\bcite{125}中有更详细讨论。
\end{note}
